% !TeX program = xelatex
\documentclass[11pt,a4paper]{article}
\usepackage[margin=23mm]{geometry}
\usepackage{fontspec,polyglossia}
\setdefaultlanguage{russian}\setotherlanguage{english}
\setmainfont{DejaVu Sans}\newfontfamily\cyrillicfont{DejaVu Sans}
\usepackage{amsmath,amssymb,graphicx,xcolor,hyperref,enumitem}
\definecolor{NESblue}{RGB}{5,67,114}
\hypersetup{colorlinks=true,linkcolor=NESblue,urlcolor=NESblue,pdftitle={Microeconomics 3: Instructor notes},pdfauthor={Arseniy Stolyarov}}
\setlength{\parindent}{0pt}\setlength{\parskip}{5pt}\setlength{\emergencystretch}{3em}\raggedright
\title{Микроэкономика · Лекция 3\\Двойственность\\и уравнение Слуцкого\\Заметки преподавателю}
\author{Арсений Столяров}\date{10 ноября 2026}
\newcommand{\R}{\mathbb R}
\begin{document}\maketitle
\section*{Как читать эту версию}
Английская презентация следует концу присланной lecture 11 (страницы 176--180) и lecture 12 (181--198), без повторного полного разбора Кун--Таккера и Роя. Короткие доказательства разбиты на клики; вопросы можно оставлять открытыми до ответа группы. В handout виден итоговый вид каждого слайда. Эти заметки содержат ответы и уточнения, включая условия на область и производные.

\section*{Две пары}
Первая: минимизация расходов и её существование; базовые свойства спроса; задачи с решением на доске; двойственность; Шепард; компенсированный закон спроса и симметрия производных. Вторая: точный рисунок изменения цены; вывод Слуцкого; примеры Кобба--Дугласа, комплементов и квазилинейных режимов; пример Гиффена; восстановление и проверка матрицы из обычного спроса. После двойственности, после классификации доходных эффектов и после примера Гиффена стоят вопросы на понимание; ответы находятся в этих заметках. Приложение начинается с квазилинейности, затем содержит все матричные расширения, отложенные свойства расходов, подробности Шепарда и прежний резерв: CV/EV, два типа компенсации, интегрируемость, три товара и непрерывность.

\section*{Что не следует утверждать}
Собственный компенсированный эффект неположителен, но не обязательно строго отрицателен. Вогнутость расходов не требует выпуклых предпочтений и не бывает строгой на ценовых лучах. Равенство множеств спроса при соответствующих параметрах не даёт равенства их частных производных. На изломах и границах режимов проверяем существование производных отдельно. Тесты матрицы — необходимые ограничения; глобальную теорему интегрируемости без условий на область не заявляем.

\section*{Послайдовые пояснения и ответы}

\subsection*{1. Microeconomics: Lecture 3}
Продолжаем ровно после Роя. Если минимизацию расходов уже успели посмотреть в резерве второй лекции, первые определения пройдите быстрее. Кун—Таккера заново не выводим.

\subsection*{2. What we know so far}
Коротко напомните первые две лекции: предпочтения и их представление, бюджет и выбор, маршаллианский спрос и углы, Кун--Таккер, непрямая полезность и огибающая. Формулу Роя отдельно не повторяем.

\subsection*{3. What we will learn today}
Маршрут: минимизация расходов и геометрия компактного отсечения; задачи и двойственность; Шепард, компенсированный закон спроса, симметрия и скалярное тождество взвешенных производных. После перерыва — история Слуцкого, вывод разложения, вопросы о нормальности/инфериорности, примеры, история и пример Гиффена. На основной части больше нет матричного блока. Квазилинейность первая в приложении, только если останется время.

\subsection*{4. What happens when prices change?}
Начните с вопроса о том, что меняется при росте цены. Первый механизм — относительная дороговизна товаров. Второй — изменение покупательной способности при прежнем денежном доходе. Затем сформулируйте потребность отделить один механизм от другого. Компенсированный спрос вводим как инструмент такого разделения.

\subsection*{5. The expenditure minimisation problem}
Сначала цель, потом ограничения, потом имена объектов. Здесь не фиксируем бюджет. H — набор самых дешёвых способов достижения цели; e — общая минимальная стоимость. При недостижимой цели множество пусто, поэтому достижимость надо оговорить.

\subsection*{6. Reach a target indifference curve at minimum cost}
Пять состояний: сначала одна кривая безразличия; затем область не хуже целевого уровня; потом параллельная линия больших расходов справа от оптимума; затем слишком дешёвая линия слева; наконец минимальная допустимая линия и точка касания. Числовой функции полезности и делений осей на слайде нет. Технически геометрия построена по той же гладкой кривой, чтобы пример был внутренне согласован; её формула студентам не нужна. На этом рисунке предполагается обычная монотонная геометрия предпочтений. Касание здесь иллюстрация, а не универсальное необходимое условие для изломов и границ.

\subsection*{7. The assumptions we will use}
Эти базовые условия действуют дальше, если не сказано иное. Для существования и многих свойств e локальная ненасыщаемость не нужна. Для обратимости e и v она нужна. Все заявления о дифференцировании локальные: в окрестности рассматриваемого параметра. Для совместной непрерывности по цели удобно брать внутреннюю достижимую цель выше u(0); на нижней границе речь об относительной непрерывности.

\subsection*{8. Does a minimum exist?}
Сначала утверждение, потом вопрос. Студенты должны найти замену компактности всего допустимого множества. Достаточно компактного подмножества, вне которого все варианты заведомо дороже уже имеющегося допустимого.

\subsection*{9. Cut away the expensive bundles}
Промежуточный рисунок между утверждением о существовании и доказательством. Четыре состояния: неограниченное верхнее контурное множество; выбранная допустимая точка; ограничение расходов по стоимости этой точки; выделенное пересечение K. На картинке за кадром используются кривая $x_1x_2=4$, цены (1,1) и точка (3,2) стоимостью 5. Пересечение по горизонтали идёт от 1 до 4. Его границы включены. Это компакт, хотя исходное множество допустимых точек неограниченно. Все отсечённые допустимые наборы дороже уже выбранного, так что минимум не теряется. Студентам численная функция не нужна.

\subsection*{10. Existence: cut away expensive bundles}
Непрерывность u даёт замкнутость верхнего контурного множества. Ограничение на расходы и строго положительные цены дают покоординатные верхние границы. Важно, что все исключённые допустимые точки строго дороже выбранной точки, которая осталась в K.

\subsection*{11. No excess utility}
Это полезный короткий результат: непрерывности достаточно, LNS не требуется. Особый случай h=0 надо проверить отдельно, иначе строгое удешевление неверно. Если целевой уровень ниже u(0), нулевой набор мог бы оптимально давать больше цели; поэтому нижняя граница в условии существенна.

\subsection*{12. Change the currency, not the chosen bundle}
Не путать с маршаллианской однородностью: там одновременно масштабируем цены и доход. Здесь доход не является параметром и масштабируем только цены. Линейное масштабирование всех расходов меняет значение минимума, но не его аргументы.

\subsection*{13. Convex preferences and Hicksian choices}
Существование уже есть; строго выпуклые предпочтения дают не более одного оптимума, значит ровно один. Для первого вывода используем выпуклость верхнего контурного множества и линейность стоимости. Для второго оба минимума лежат на одной кривой по предыдущему утверждению. Кривизна конкретной функции полезности здесь не нужна.

\subsection*{14. Buy a Cobb--Douglas standard of living}
На слайде только задача: найти общий хиксианский спрос и расходы для корня из произведения, затем посчитать при ценах (4,1) и цели 3. Остановитесь и решите на доске. Следующий слайд содержит ответ с постепенным появлением. Положительная цель исключает оси.

\subsection*{15. Cobb--Douglas: the dual solution}
Первая страница только напоминает цель. Затем по кликам появляются два уравнения, формулы спроса и расходов, численный ответ. На доске можно подставить $x_2=\bar u^2/x_1$ и минимизировать $p_1x_1+p_2\bar u^2/x_1$. Производная даёт равенство расходов $p_1h_1=p_2h_2$. При ценах (4,1) и цели 3 получаем (1.5,6), расходы 12.

\subsection*{16. Perfect complements: find the cheapest bundle}
Сначала только полезность, цены, цель и просьба найти решение. После решения на доске раскрывайте: цель равносильна двум нижним границам; минимизатор равен (ubar,ubar), стоимость (p1+p2)ubar. Ни касание, ни производные самой min-функции не нужны. Последний клик напоминает, что излом прямой полезности совместим с гладкой функцией расходов.

\subsection*{17. Perfect substitutes: which route is cheapest?}
Сначала дана сама задача минимизации расходов без графика и без ответа. Затем попросите найти все минимизаторы, включая равные цены. Разберите на доске, после чего открывайте стоимость, строгие ценовые режимы и весь отрезок минимизаторов при равенстве цен. Излом e можно обсудить словами; график удалён.

\subsection*{18. Two problems, the same boundary}
Это центральная теорема первой половины. Утверждаем равенство множеств, не только специально выбранных оптимумов. Цель выше u(0) обеспечивает положительную соответствующую сумму расходов; нижнюю границу можно рассмотреть отдельно. В доказательствах не нужна выпуклость предпочтений.

\subsection*{19. Proof: the cheapest way to buy the maximum}
Если h давал больше максимума, это уже противоречие определению v. Если ровно максимум, он оптимален с непотраченными деньгами, что запрещено LNS. Это хороший момент спросить, где именно понадобилась локальная ненасыщаемость.

\subsection*{20. Proof: the most utility the minimum cost buys}
Почему z тоже минимизатор? Он достигает цели и стоит не больше её минимальной стоимости. Поэтому не может стоить меньше и стоит ровно столько. Тогда предыдущая лемма запрещает ему давать избыточную полезность. В этом направлении основной инструмент — непрерывность и отсутствие избыточной полезности.

\subsection*{21. Proof: match the sets, not just their values}
Оба множества записываются через пересечение верхнего контурного множества и бюджетного. Знаки ≥ и ≤ допустимы: само определение максимума/минимума исключает строгое улучшение. Этот аргумент аккуратно обходится с множественным спросом.

\subsection*{22. Check duality for Cobb--Douglas}
На слайде только задача. Ответ для преподавателя: $v=m/(2\sqrt{p_1p_2})$, $e=2\bar u\sqrt{p_1p_2}$, $x_i=m/(2p_i)$, $h_1=\bar u\sqrt{p_2/p_1}$ и $h_2=\bar u\sqrt{p_1/p_2}$. Подстановка даёт $e(p,v(p,m))=m$, $v(p,e(p,\bar u))=\bar u$ и совпадение спроса при соответствующих параметрах. При дифференцировании позже важно помнить, какой параметр фиксируется.

\subsection*{23. Checkpoint}
Три вопроса, ответы только здесь:
\begin{enumerate}
\item При удвоении всех цен хиксианские количества не меняются, расходы удваиваются.
\item Новая цель $2\cdot7+10=24$. Это та же кривая безразличия с новой числовой меткой. Старую метку 7 оставлять нельзя: она соответствовала бы другому уровню исходной полезности.
\item Нет при условиях теоремы двойственности. Для $\bar u=v(p,m)$ имеем $X(p,m)=H(p,\bar u)$, поэтому два маршаллианских оптимума остаются двумя хиксианскими. Речь о соответствующей цели, а не о произвольном другом уровне полезности.
\end{enumerate}
Четвёртый вопрос о вогнутости e удалён, поскольку само свойство теперь находится в приложении.

\subsection*{24. Shephard’s lemma}
Правильное написание — Shephard. При наших условиях непрерывности u, положительных ценах и достижимой фиксированной цели оптимумы локально ограничены; единственность даёт дифференцируемость через теорему об огибающей для минимума. Дифференцируемость самой h не предполагается. Прямое доказательство следующего слайда работает даже без гладкости u.

\subsection*{25. Recover quantities from cost}
Дайте студентам сделать производные прежде клика. Здесь достаточно знания e, функцию полезности заново оптимизировать не надо. Для субститутов при равных ценах производной нет; в строгом режиме одна производная равна цели, другая нулю.

\subsection*{26. The compensated law of demand}
Верхние индексы 1 и 2 обозначают две ситуации с разными векторами цен; нижний индекс i — товар. Пишем два выбранных минимизатора как $h^1$ и $h^2$, вместо смешения букв h и k. Одинаковая цель делает оба набора допустимыми в обеих задачах. Если меняется только цена i, скалярное неравенство превращается в $(p_i^2-p_i^1)(h_i^2-h_i^1)\le0$. Это не требует гладкости, выпуклости или единственности.

\subsection*{27. Two cost comparisons prove the law}
Два сравнения с одинаковой целью: первый набор дешевле второго при первых ценах, второй дешевле первого при вторых. Складываем и группируем. Именно эта одинаковая цель позволяет использовать чужой набор как допустимый. Индексы ситуаций сохранены в точности как на предыдущем слайде.

\subsection*{28. Symmetry of compensated price effects}
В основной части оставлено равенство перекрёстных производных хиксианского спроса. Условие: e дважды непрерывно дифференцируема по ценам в окрестности. Шепард переводит количества в первые производные e, после чего используем равенство смешанных производных. На следующем слайде сформулируем скалярное тождество взвешенных производных. Матричную запись и отрицательную полуопределённость оставляем в приложении для семинара.

\subsection*{29. Substitution has zero first-order cost}
Важно уточнить исходную формулировку: нулю равно скалярное произведение вектора цен на вектор изменения компенсированного спроса, а не каждое отдельное произведение. При фиксированной полезности $\sum_jp_j\,\partial h_j/\partial p_i=0$. Доказываем дифференцированием тождества минимальной стоимости $e=p\cdot h$ и леммой Шепарда. Благодаря ранее доказанной симметрии можно также написать $\sum_jp_j\,\partial h_i/\partial p_j=0$; это же следует из однородности h нулевой степени. Первое выражение суммирует изменения разных количеств при изменении одной цены; второе меняет все цены пропорционально. Никаких матриц на этом слайде нет. Нулевые расходы от изменения количеств не означают нулевого изменения общей стоимости при росте цен: прямой эффект $de=h\cdot dp$ остаётся.

\subsection*{30. Substitution and income effects}
На слайде нет функции полезности, численных цен, дохода или делений осей. Линии появляются по одной: исходный бюджет; исходная кривая и A; новый бюджет после роста первой цены; новая кривая и C; компенсирующий бюджет и B. Последними отдельно появляются стрелки A к B и B к C. Первая сохраняет исходную полезность, вторая новые цены. Геометрия согласована с прежним примером Кобба--Дугласа, но никаких вычислений для понимания рисунка не требуется. Это конечное хиксианское разложение; дифференциальный Слуцкий следует дальше.

\subsection*{31. The choices match. Their price derivatives do not.}
Перед формальным выводом попросите студентов назвать явное и косвенное влияние цены. Это тот же chain rule, который уже встречался в огибающей. Здесь слагаемое с производной дохода не исчезает: меняется именно аргумент m=e(p,ubar).

\subsection*{32. Eugen Slutsky and the consumer’s budget}
Евгений Евгеньевич Слуцкий (1880--1948) работал в экономике, математике и статистике. В 1915 году опубликовал на итальянском статью «Sulla teoria del bilancio del consumatore» в Giornale degli Economisti: «К теории бюджета потребителя». Здесь достаточно минуты исторического контекста, затем переходим к нашему выводу через двойственность и правило цепочки. Преподаватель прислал фотографию в чат. Для включения нужен оригинальный файл assets/slutsky-portrait.jpg; когда он доступен, сборка вставляет его автоматически. В текущей облачной папке оригинала пока нет; сгенерированное изображение удалено.

\subsection*{33. Slutsky: the chain rule does the work}
Оценки аргументов важны: производные x берутся при доходе e(p,ubar), производные h — при фиксированном ubar. После подстановки ubar=v(p,m) это одна исходная точка. Не дифференцируем v(p,m) в последней подстановке: производные уже взяты при фиксированной цели.

\subsection*{34. The Slutsky equation}
Ловушка с индексами: подорожавший товар j определяет потерю покупательной способности; товар i — тот, чей спрос меняется. Параметры p и m справа — исходные. Эффект дохода от роста цены равен $-x_j\,\partial x_i/\partial m$. Нельзя называть само положительное произведение эффектом без оговорки знака.

\subsection*{35. When must ordinary demand slope down?}
Используем локальные определения. При положительном потреблении Гиффен обязательно низший, но низший товар не обязательно Гиффен. Для нормального товара вывод нестрогий: возможны нулевые производные. Товары не имеют этих ярлыков навсегда — свойство может зависеть от дохода и цен.

\subsection*{36. Checkpoint: normal and inferior goods}
Ответы: F, T, T, T. Во всех вопросах фиксированы цены и рассматриваются локальные доходные производные. Дифференцируем исчерпание бюджета:
\[\sum_i p_i x_i(p,m)=m\quad\Rightarrow\quad\sum_i p_i\frac{\partial x_i}{\partial m}=1.\]
При двух товарах обе производные не могут быть отрицательны: положительно взвешенная сумма должна равняться 1. Оба товара могут быть строго нормальными, например в Коббе--Дугласе. Один нормальный и один инфериорный возможны в нашем примере Гиффена: доходные производные $-1/(2p_2-p_1)$ и $2/(2p_2-p_1)$.

При трёх товарах два могут быть инфериорными, если третий поглощает дополнительный доход. Явный локальный пример: положим $q=\min\{x_1,x_2\}$ и $u=\min\{q+x_3,3q+2x_3-10\}$. При ценах $(0.6,0.6,1)$ и доходе 8 оптимум $(2.5,2.5,5)$. Цена композита q равна 1.2, и это уже решённый двухтоварный пример. Локальные доходные производные равны $(-1.25,-1.25,2.5)$, а их взвешенная сумма $0.6(-1.25)+0.6(-1.25)+2.5=1$. Дополнительные количества одного из первых двух товаров сверх другого бесполезны и дороги, поэтому оптимальны равные количества. На слайде оставлены только утверждения; ответы проговорите после попытки группы.

\subsection*{37. Cobb--Douglas: split an own-price effect}
Сначала дифференцируем h при фиксированной цели, затем подставляем ubar=m/(2√p1p2). Если подставить цель раньше и дифференцировать итог m/(2p1), получим маршаллианскую производную, а не хиксианскую. Это важнейшая техническая ловушка этой лекции.

\subsection*{38. No total effect does not mean no substitution}
Вопрос уточняем: при фиксированном денежном доходе итогового перекрёстного эффекта нет, но при компенсации он положителен. Это объясняет вертикальное расположение A и C на рисунке: выросший спрос на второй товар при компенсации исчезает из-за снижения реального дохода.

\subsection*{39. Perfect complements: all income, no substitution}
Этот пример полезен сразу по трём причинам: прямая полезность негладкая, функции спроса и расходов гладкие, собственный компенсированный эффект не обязан быть строго отрицательным. Два товара — валовые комплементы, но хиксианские перекрёстные эффекты равны нулю.

\subsection*{40. Can a higher price make you buy more?}
Сначала расскажем картофельную историю, затем проверим механизм на явной функции полезности. Рост спроса на престижный товар из-за сигнала качества не является чистым сравнением по собственной цене в неизменных предпочтениях над теми же физическими товарами. После истории дадим полностью проверяемый искусственный пример. Он вогнутый как минимум двух аффинных функций, значит предпочтения выпуклы, и строго монотонный.

\subsection*{41. Robert Giffen and the potato story}
Роберт Гиффен (1837--1910) — шотландский экономист и статистик. Маршалл в Principles of Economics связывал с Гиффеном пример спроса на хлеб у бедных семей. Традиционный рассказ про картофель — учебная версия механизма; его не следует выдавать за доказанный эмпирический факт об Ирландском голоде.

Семья живёт в основном на картофеле, но иногда покупает мясо. Картофель даёт самые дешёвые калории и занимает большую долю бюджета. Он дорожает при неизменных денежном доходе и цене мяса. Семья сокращает мясо и заменяет потерянные калории картофелем, который всё ещё дешевле на калорию. Его покупка может вырасти. Изменение относительных цен само по себе толкает от картофеля, но эффект снижения реального дохода может оказаться сильнее. При компенсации полезности положительный собственный эффект невозможен.

Слева краткая историческая справка, справа по кликам рассказ. Для присланного преподавателем гравированного портрета подготовлен путь assets/giffen-portrait.png. Оригинала в облачной папке пока нет; сгенерированное изображение удалено. При появлении оригинала сборка вставляет его автоматически. Формальный пример дальше независимо подтверждает математическую возможность эффекта при выпуклых предпочтениях.

\subsection*{42. A subsistence kink generates Giffen demand}
Полосы условий обеспечивают внутренний оптимум на изломе. Между наклонами градиентов (1,1) и (3,2) отношение цен лежит строго между 1 и 1.5; бюджет имеет единственный максимум в точке переключения. Более строго: вес eta= (p1−p2)/(2p2−p1) лежит (0,1), а взвешенный градиент (1+2eta,1+eta) пропорционален p. Активны обе ветви, поэтому их равенство и бюджет определяют единственную точку. Отрицательные значения полезности вне рассматриваемой области допустимы: шкала произвольна. Для интуиции это два линейных ограничения на уровень жизни, одно содержит фиксированный порог.

\subsection*{43. The income effect wins}
Числа точные: при 1.2 знаменатель .8, при 1.3 .7. Исходная полезность 7.5. При фиксированной цели 7.5 оба линейных ограничения связывают h=(2.5,5) при небольших изменениях цен, так что компенсированный эффект ноль. Рост обычного спроса полностью объясняется доходным эффектом. При новой цене полезность 50/7≈7.143 ниже исходной. Это не противоречит компенсированному закону спроса. Общая локальная формула h=(10−ubar,2ubar−10), e=p1(10−ubar)+p2(2ubar−10), пока обе координаты положительны и отношение цен в указанной полосе.

\subsection*{44. Checkpoint}
Нижние индексы после запятой обозначают частные производные. Ответы:
\begin{enumerate}
\item Нет: $x_{1,p_1}=-3-2(-1)=-1<0$. Товар инфериорный, но не Гиффена.
\item Нет: в Коббе--Дугласе компенсированный перекрёстный эффект и эффект дохода точно компенсируются.
\item Нет. При фиксированном целевом уровне действует компенсированный закон спроса, даже без дифференцируемости.
\item Нет. Уже рассмотренный пример Гиффена имеет вогнутую полезность как минимум двух аффинных функций, значит выпуклые предпочтения. Гиффен не свидетельствует о нарушении рациональности или выпуклости.
\end{enumerate}
Вопрос о квазилинейности заменён четвёртым вопросом о Гиффене, поскольку сам разбор квазилинейности теперь первый в приложении.

\subsection*{45. Substitutes and complements: specify the experiment}
Не употребляем слово substitutes без уточнения, если обсуждаем производные. Обычно в названии perfect complements речь о форме предпочтений; нулевой хиксианский перекрёстный эффект этому не противоречит. Нулевая производная — пограничный случай, не строгая подстановка или дополняемость. Симметрия относится к изменениям количеств на единицу цены, не к эластичностям.

\subsection*{46. Checkpoint}
Ответы: F, F, F, T. Первое неверно уже из-за однородности e первой степени: на любом ценовом луче линейность исключает строгую вогнутость. Второе: Leontief даёт нулевые производные. Третье: у смещённого произведения (1+x1)(1+x2) во внутреннем режиме xi=(m+p1+p2)/(2pi)−1; перекрёстные производные 1/(2p1) и 1/(2p2) неодинаковы при разных ценах. Четвёртое следует из собственного уравнения Слуцкого при существовании производных и xi>0.

\subsection*{47. One structure, several ways to use it}
Верните лекцию к четырём объектам и их назначениям. Вопрос не только в решении примеров оптимизации: структура даёт ограничения на данные и способ оценивать компенсацию. Нельзя читать двустороннюю стрелку как восстановление уникальной кардинальной полезности: идентификация только ординальная, и при неполных данных не единственная.

\subsection*{48. Read next}
Никакого нового материала третьей лекции задним числом не добавляем в опубликованную домашку. Её дедлайн 14 ноября, темы — первые две лекции. Про будущую домашку пока нет новых объявлений. Ссылки на загруженные лекции не нужны на студенческом слайде; сведения об источниках находятся в SOURCE\_NOTES.md.

\subsection*{49. Quasilinearity: income effects depend on the regime}
Первый необязательный слайд сразу после конца основной лекции. Его можно разобрать, только если останется время.

Это ровно квазилинейная функция из второй лекции, полезное возвращение к режимам. В строгом all-q режиме хиксианское количество q задаётся целью, пока цена остаётся в том же режиме, поэтому производная по r ноль. На границе m=4−r спрос не дифференцируем: обычную гладкую формулу там не заявляем. При r>4 q=0 локально, все эти производные q равны нулю. Цена numeraire нормирована к единице.

\subsection*{50. What can only increase expenditure?}
По ценам нестрогое возрастание: потребление подорожавшего товара может быть нулевым. По достижимой цели выше нижней границы — строгое. Во втором доказательстве сначала вложенность допустимых множеств даёт нестрогое неравенство; затем исключаем равенство. LNS для этих сравнений не требуется.

\subsection*{51. Expenditure is concave in prices}
Здесь часто автоматически пишут неправильный знак. Функция расходов — нижняя огибающая линейных по ценам функций. Квазивогнутость полезности и выпуклость предпочтений не используются. Не утверждайте строгую вогнутость: уже однородность первой степени исключает её вдоль ценовых лучей.

\subsection*{52. Concavity: one bundle, two price vectors}
Доказательство на три клика. Не складываем различные оптимальные наборы и не применяем выпуклость допустимого множества. Берём один оптимальный набор при смешанных ценах и сравниваем его стоимость с минимумами в двух других задачах.

\subsection*{53. Utility labels can change; the cost cannot}
Смена шкалы требует сменить численный целевой уровень. Нельзя заменить u на логарифм и оставить прежнее число цели. Утверждение о непрерывности по обоим параметрам использует LNS: можно слегка улучшить минимизатор, чтобы он оставался допустимым для немного большей цели. Полное доказательство на резервном слайде; по ценам при фиксированной цели непрерывность следует также из конечной вогнутости на открытом положительном ортанте.

\subsection*{54. Why the quantity appears in the derivative}
Для верхней оценки h допустим при новых ценах; для нижней $h_t$ допустим при старых. При t<0 порядок после деления меняется, но предел тот же. Локальная ограниченность: используем один фиксированный допустимый набор для верхней границы затрат, а цены в окрестности отделены от нуля. Любая предельная точка оптимумов оптимальна при исходных ценах; единственность заставляет её совпасть с h. Аналогичный аргумент по любому направлению даёт полную дифференцируемость; координатный слайд показывает интуицию.

\subsection*{55. What if demand is not unique?}
Доказательство второй строки: подставьте q=p+td для произвольного направления d и обоих знаков t. Получаем градиент e, скалярно умноженный на d, равен h·d для всех d, следовательно h=градиент. Каждый минимизатор даёт верхнюю опорную плоскость вогнутой e. Поэтому различные минимизаторы исключают дифференцируемость e.

\subsection*{56. The matrix behind compensated substitution}
Фиксируем конвенцию: строка i — количество товара i, столбец j — цена j. S в точке соответствует именно компенсированным производным, а не обычному якобиану спроса. C² e в окрестности даёт равенство смешанных производных. Отрицательная полуопределённость сильнее условия неположительности одной диагонали.

\subsection*{57. Why symmetry and the quadratic inequality?}
Произвольный вектор z — направление одновременного изменения цен. Положительность исходных цен позволяет взять маленькую окрестность t=0 для любого z. Матрица не обязана быть отрицательно определённой: однородность сейчас даст обязательное нулевое направление.

\subsection*{58. A common price increase changes no compensated quantities}
Здесь лучше механическое доказательство через однородность, чем рисунок касательной: оно не требует внутреннего оптимума. Приведённая в исходнике геометрическая интерпретация менее универсальна на границе. Поскольку p ненулевой, S всегда вырождена. Нельзя требовать строгой отрицательной определённости на всём пространстве.

\subsection*{59. A surprising restriction with only two goods}
Матричное расширение для приложения и семинара.

Это следствие двух уже доказанных свойств матрицы, а не новой геометрической догадки. Leontief даёт ноль. Наличие третьего товара позволяет компенсировать рост цен изменением третьего количества, поэтому отрицательные перекрёстные эффекты не запрещены. В резерве есть явный пример с логарифмом комплементарной пары и третьим товаром: $u=\log(\min\{x_1,x_2\})+x_3$ на указанной там области.

\subsection*{60. Construct compensated effects from ordinary demand}
Матричное расширение для приложения и семинара.

$D_m x$ — столбец, x транспонированный — строка, поэтому произведение внешнее. Ошибка ориентации меняет матрицу. Наблюдаемые оценки должны относиться к одному стабильному потребителю или корректно моделируемой неоднородности. Не заявляем, что произвольная средняя по разным людям автоматически имеет ту же структуру. Глобальная достаточность тестов требует отдельной теоремы интегрируемости и условий на область.

\subsection*{61. Cobb--Douglas passes the matrix test}
Матричное расширение для приложения и семинара.

Обычный якобиан здесь диагонален: Dpx имеет −m/(2pi²) на диагонали и нули вне. Доходная поправка добавляет положительные элементы, давая S. Квадрат немедленно проверяет полуопределённость; его нулевое направление z пропорционально p — общий рост всех цен.

\subsection*{62. Could these choices come from our consumer model?}
Матричное расширение для приложения и семинара.

Не сразу открываем ответ. На точке Dpx=[[0,−1],[−1,0]]. Доходная поправка — матрица всех элементов .5. Положительные количества, Вальрас и отсутствие денежной иллюзии сами по себе не гарантируют рационализируемости.

\subsection*{63. A symmetric matrix can still fail}
Матричное расширение для приложения и семинара.

Собственные значения 0 и 1. Здесь особенно наглядно, что симметрия и нулевое направление не заменяют отрицательную полуопределённость. Это контрпример с обычными нормальными товарами по доходу, но не с допустимой компенсированной реакцией.

\subsection*{64. Reserve: compensation and equivalent variation}
Соглашение о знаках явно написано: положительное число — денежная оценка потери при росте цен. В литературе встречаются противоположные знаки. Эти две величины обычно различаются. Они измеряют денежную компенсацию индивидуального изменения, а не универсальную кардинальную шкалу благосостояния разных людей.

\subsection*{65. Reserve: how much compensation?}
Это та же картинка A,B,C. Новая сумма 8+CV=8√2 финансирует B и восстанавливает полезность 4. При старых ценах доход 8−EV=4√2 даёт полезность 2√2. Можно проверить численно без новой оптимизации.

\subsection*{66. Reserve: two meanings of compensation}
Название уравнения Слуцкого не означает, что на предыдущей картинке компенсация была по стоимости старой корзины. Рисунок использовал хиксианскую компенсацию. Для конечных изменений новая бюджетная линия через старую корзину обычно позволяет более высокий уровень полезности, чем исходный. В примере 12≥8√2. Дифференциальное разложение совпадает в первом порядке.

\subsection*{67. Reserve: can we work backwards from demand?}
Не обещаем, что четыре условия в любой произвольной области автоматически дают глобально определённую единственную полезность. Нужны стандартные условия гладкости, области и продолжения решений. Симметрия — условие согласованности системы дифференциальных уравнений, полуопределённость связана с вогнутостью восстановленной e. Даже при успехе восстанавливается только ординальный представитель; поведения на никогда не выбираемых наборах может быть недостаточно для идентификации исходных предпочтений.

\subsection*{68. Reserve: genuine net complements need another good}
Этот резервный пример отдельно меняет область: x1,x2 строго положительны из-за логарифма. Не применяем к нему без оговорки утверждения с u(0). Для минимизации полезно положить x1=x2=q, иначе избыточное количество одного товара только дорожает. При активной цели x3=ubar−log q, q≤exp(ubar). Стоимость cq+p3(ubar−log q) имеет уникальный минимум q=p3/c, если он строго ниже exp(ubar); это и есть условие режима. Производная h1 по p2 отрицательна. Симметричная производная h2 по p1 такая же. Компенсация повышает x3; вычисление на конкретном режиме достаточно для иллюстрации. Существование тут следует из явного одномерного решения, а не из непрерывности u на замкнутом ортанте.

\subsection*{69. Reserve: continuity of expenditure}
В нижней оценке выбираем подпоследовательность значений, сходящуюся к liminf, и затем сходящуюся подпоследовательность оптимальных наборов. Цены в окрестности положительного p отделены от нуля, поэтому ограниченные расходы дают ограниченные количества. В верхней оценке сначала выбираем одну слегка улучшенную точку z, потом переходим к пределу параметров, потом приближаем z к h. Не меняем порядок этих шагов. На граничной цели u(0) нужна относительная формулировка.

\end{document}
